% TOP_PROBLEM: {"id": 3212, "title": "Graceful Tree Conjecture", "category_id": 3, "category": {"id": 3, "name": "graph_theory", "display_name": "Graph Theory", "description": "Problems involving graphs, networks, and their properties.", "slug": "graph-theory", "order_index": 3, "created_at": "2026-07-31T15:26:25.671Z"}, "status": "open", "problem_number": "OPG-439", "set_id": 12, "statusQualification": "Injectivity of vertex labels is explicit, correcting the omission noted in the source discussion."}
% FIRST_POSTED: 2026-10-01
% ENTRY_KIND: statement_only
% UnsolvedMath ID 3212; source catalogue code OPG-439
% !TeX program = lualatex
% Definition and sources only; this file is not an LLM solution attempt.
\documentclass[11pt,a4paper]{article}
\usepackage[margin=25mm]{geometry}
\usepackage{fontspec}
\setmainfont{lmroman10-regular.otf}[BoldFont=lmroman10-bold.otf,
 ItalicFont=lmroman10-italic.otf,BoldItalicFont=lmroman10-bolditalic.otf]
\usepackage{amsmath,amssymb,mathtools}
\usepackage{unicode-math}
\setmathfont{latinmodern-math.otf}
\AtBeginDocument{\let\setminus\smallsetminus}
\usepackage{xurl,hyperref}
\hypersetup{colorlinks=true,urlcolor=blue}
\setcounter{secnumdepth}{0}
\setlength{\parindent}{0pt}
\setlength{\parskip}{5pt}
\setlength{\emergencystretch}{3em}
\begin{document}
\section*{Graceful Tree Conjecture}\label{3212}
{\small UnsolvedMath ID 3212 · OPG-439 · Graph Theory · OpenGarden}
\subsection{Definitions and mathematical statement}
Let $T=(V,E)$ be a finite tree, meaning
a nonempty simple connected graph with
no cycle. Write $m=|E|$, so $|V|=m+1$.
A graceful labelling is an injective
function
\[
 f:V\longrightarrow\{0,1,\ldots,m\}
\]
for which the induced edge labels
$|f(u)-f(v)|$, $uv\in E$, are pairwise
distinct. Since there are $m$ edges
and each edge difference lies in
$\{1,\ldots,m\}$, this is equivalent
to requiring
\[
 \{|f(u)-f(v)|:uv\in E\}=\{1,\ldots,m\}.
\]
For a tree the vertex labelling is
necessarily a bijection, since domain
and codomain have the same size.

The Graceful Tree Conjecture states
that every finite tree admits a
graceful labelling. Both the vertex
labels and the absolute edge differences
must satisfy the specified conditions
simultaneously. Repeated vertex labels
are forbidden even if the edge
differences happen to be distinct.

The one-vertex tree, with $m=0$,
has the label zero and no edges,
so it satisfies the definition
vacuously. In general, no root,
orientation, or ordering of the
vertices is prescribed; the labelling
may be chosen freely to suit the
tree. Enlarging the label interval
or allowing fewer than $m$ distinct
edge differences would give weaker
labelling problems. The target here
requires exactly the interval from
zero to $m$ and every positive
difference in that interval once.
\subsection{Short English statement}
Can every tree be numbered from zero to its number of edges so the differences across its edges are exactly all positive labels, without repetition?
\subsection{Sources}
\begin{itemize}
\item[S1] ULAM.AI and the UnsolvedMath Contributors, supplied problems.json, record 3212 (OPG-439), OpenGarden collection. Catalogue identity and source transcription; CC BY 4.0.
\url{https://huggingface.co/datasets/ulamai/UnsolvedMath}.
\item[S2] Open Problem Garden, Graceful Tree Conjecture. Original problem and discussion; the mathematical formulation below is expanded from the supplied source record.
\url{http://www.openproblemgarden.org/op/graceful_tree_conjecture}.
\end{itemize}
\end{document}
